\section{DFA Transition Function}

The DFA’s transition function, \(\delta\), is defined as follows:
\begin{align*}
{%- for state, transitions in tf.items() -%}
    {%- for symbol, target in transitions.items() -%}
        \delta({{ state }},\,{{ symbol }}) &= {{ target }} {%- if not loop.last %}\\[0.5em]{%- endif %}
    {%- endfor %}
{%- if not loop.last %}\\[1em]{%- endif %}
{%- endfor %}
\end{align*}

\section{Indistinguishability Table Steps}

We apply the table‑filling (indistinguishability) algorithm iteratively.  For each iteration, we show the triangular table and provide a short explanation of which pairs were marked.

{%- for idx in range(indis_plots|length) %}
\subsection*{Table Step {{ idx+1 }}}

\textbf{Explanation:} {{ indis_details[idx].description }}

\begin{center}
    \includegraphics[width=0.6\linewidth]{plots/indis_step{{ idx+1 }}.png}
\end{center}
\vspace{1em}
{%- endfor %}

\section{Minimized DFA}

After no further distinguishable pairs can be marked, the remaining unmarked state‑pairs are equivalent and are merged.  The resulting minimized DFA is:

\begin{center}
    \includegraphics[width=0.7\linewidth]{minimized_dfa.png}
\end{center}
